\section{Discussion}
\label{sec:discussion}

\subsection{Implications of Lunar Magnetic Heterogeneity}
The central finding of this analysis is that while the lunar crustal magnetic field exhibits pronounced spatial heterogeneity across the lunar globe, all evaluated candidate and historical landing sites represent severe hypomagnetic field (HMF) environments from a biological perspective. Even the most intense lunar magnetic anomalies, which reach $\sim\SI{20}{\nano\tesla}$ at \SI{30}{\kilo\meter} altitude, are more than three orders of magnitude ($>2,500\times$) weaker than Earth's geomagnetic field (\SI{25}{\micro\tesla} to \SI{65}{\micro\tesla}). At candidate Artemis III South Polar sites ($< \SI{0.8}{\nano\tesla}$), the local crustal magnetic field is more than five orders of magnitude ($>60,000\times$) weaker than Earth's GMF. Consequently, any biological system deployed to the lunar surface will experience a profound, sustained departure from its evolutionary magnetic baseline.

However, the relative differences between lunar sites are not trivial. A site situated within a moderate crustal anomaly ($\sim\SI{5}{\nano\tesla}$) offers localized interaction with the solar wind plasma, forming mini-magnetospheres that may partially deflect low-energy solar energetic particles \citep{Mitchell2008}. In contrast, near-null sites ($< \SI{0.5}{\nano\tesla}$) experience direct solar wind and plasma impingement without local crustal magnetic perturbation. Whether this localized shielding difference translates into altered biological responses remains a critical open scientific question.

\subsection{Biological Risks for Surface Operations}
Our synthesis of the magnetobiology literature indicates that prolonged exposure to HMFs poses multifaceted risks to biological payloads and life support systems:
\begin{enumerate}
    \item \textbf{Plant Growth and Bioregenerative Life Support (BLSS):} For crop species and model plants such as \textit{Arabidopsis thaliana}, near-null magnetic fields disrupt cryptochrome and phytochrome signaling, delaying floral transition and altering root gravitropic morphology via PIN2 auxin redistribution \citep{Agliassa2018, Narayana2018}. Furthermore, HMF-induced ROS accumulation and impaired iron uptake \citep{Maffei2014, Narayana2021} could impair crop yields in long-duration surface greenhouses.
    \item \textbf{Microbial Ecology and Pathogenicity:} Bacterial populations in HMFs exhibit altered growth kinetics, modified biofilm architectures, and shifted metabolic flux \citep{Creanga2009, Zhang2017}. Managing closed-loop habitat microbiomes will require accounting for magnetic field reductions.
    \item \textbf{Mammalian and Astronaut Health:} In vitro and rodent model studies demonstrate that HMF conditions exacerbate musculoskeletal bone demineralization induced by reduced gravity \citep{Xu2012, Mo2014}, perturb cell cycle transcriptomes \citep{Mo2012}, and modify radical pair spin-state recombination kinetics in DNA damage responses \citep{Binhi2017, Hore2016}.
\end{enumerate}

\subsection{Limitations}
Several limitations must be noted. First, the field values used in this study represent orbital determinations at an altitude of \SI{30}{\kilo\meter}. Surface magnetic field intensities directly at ground level are expected to be higher due to proximity to shallow magnetized crustal bodies (potentially reaching tens to hundreds of nanoteslas at localized anomaly crests). Second, terrestrial HMF simulation facilities typically achieve shielding down to $\sim\SI{10}{\nano\tesla}$ to \SI{100}{\nano\tesla}, which is still elevated compared to the sub-nanotesla baseline of the lunar polar craters. Third, terrestrial simulations cannot fully replicate the combined synergy of lunar hypogravity ($0.166\,g$), chronic galactic cosmic radiation (GCR), and lunar dust toxicity.

\subsection{Recommendations for Future Missions}
To address these challenges, we propose the following strategic recommendations for upcoming lunar surface exploration:
\begin{enumerate}
    \item \textbf{In Situ Surface Magnetometry:} Artemis surface payloads should prioritize calibrated vector magnetometers to record high-resolution ground-level fields at human scale.
    \item \textbf{Controlled Biological Payloads:} Surface biological experiments should integrate active Helmholtz coils or Mu-metal shielding to independently isolate magnetic field variables from radiation and hypogravity.
    \item \textbf{Magnetic Mapping for Habitat Siting:} While primary Artemis site selection focuses on illumination, slope, and volatiles, local magnetic field architecture should be factored into long-term infrastructure and biological facility planning.
\end{enumerate}
